\section*{Chapitre \ref{ch:rc:liquidity-and-funding-risk-bank-treasury} --- Risque de liquidité et de financement~; trésorerie bancaire}

\begin{solution}{exo:rc:liquidity-and-funding-risk-bank-treasury:1}
Le niveau~2A après décote vaut $0.85\times40 = 34$, plafonné à $\tfrac23\times30 = 20$ (le niveau~2 ne dépasse pas 40~\% du
stock)~; HQLA $= 30+20 = 50$.
\end{solution}

\begin{solution}{exo:rc:liquidity-and-funding-risk-bank-treasury:2}
$-6.2\times3.5~\% = -21.7~\%$ au premier ordre~; la baisse exacte est plus faible (19.2~\% pour les MBS du chapitre) parce que
de la convexité positive, et parce que la duration elle-même baisse lorsque les rendements montent.
\end{solution}

\begin{solution}{exo:rc:liquidity-and-funding-risk-bank-treasury:3}
$300\times e^{-5/4} = 86$ points de base.
\end{solution}

\begin{solution}{exo:rc:liquidity-and-funding-risk-bank-treasury:4}
Ses actifs sont plus longs que ses passifs (titres à taux fixe longs face à des dépôts échelonnés
sur cinq ans au plus)~: la banque est longue en duration, comme un receveur de taux fixe dans un swap. Elle gagne lorsque les taux baissent
et perd lorsqu'ils montent.
\end{solution}

\begin{solution}{exo:rc:liquidity-and-funding-risk-bank-treasury:5}
Les déposants non garantis perdent de l'argent si la banque fait faillite, si bien que tout doute sur sa solvabilité donne à chacun d'eux une
raison de partir les premiers~; lorsqu'ils sont peu nombreux, importants et liés entre eux, ils partent ensemble. Leur retrait est
rapide et corrélé, ce qui constitue un risque de liquidité pour la banque quel que soit son capital.
\end{solution}

\begin{solution}{exo:rc:liquidity-and-funding-risk-bank-treasury:6}
L'activité est facturée au taux de financement à terme de ses prêts (4.84~\% sur cinq ans dans l'exemple), de
sorte qu'un prêt long ne rapporte que sa marge de crédit et de service~; la prime de terme et de liquidité revient à la trésorerie,
qui le finance et le couvre. Sans \omterm{def:rc:liquidity-and-funding-risk-bank-treasury:ftp}{taux de cession interne} (funds-transfer pricing, FTP), l'activité enregistrerait la pente de la courbe comme un profit.
\end{solution}

\begin{solution}{exo:rc:liquidity-and-funding-risk-bank-treasury:7}
$+0.86$ par an avec un bêta de 0.35, $-1.87$ avec 0.80.
\end{solution}

\begin{solution}{exo:rc:liquidity-and-funding-risk-bank-treasury:8}
Détenir les titres jusqu'à l'échéance n'évite de réaliser la perte que si la banque n'a jamais besoin de vendre ou de nantir les
titres~; leur valeur économique a tout de même baissé, la banque perçoit sur eux des rendements inférieurs au marché
pendant des années, et une sortie de dépôts peut forcer une vente. Les garanties publiques suppriment le risque de crédit, et non
le risque de taux d'intérêt ni le risque de liquidité.
\end{solution}

\begin{solution}{pb:rc:liquidity-and-funding-risk-bank-treasury:1}
\textbf{1.} 17.2~milliards de dollars, soit 108~\% des fonds propres.
\textbf{2.} 3.9~milliards de dollars.
\textbf{3.} $16-17.2-3.9 = -5.1$~: négatif.
\textbf{4.} Une hausse de 200 points de base réduisait déjà l'EVE de 6.95, soit 43~\% du Tier~1, bien au-delà du seuil
de valeur aberrante de 15~\%~: le décalage de duration était visible avant toute perte.
\textbf{5.} Des swaps payeurs de taux fixe de même duration~: l'EVE devient insensible aux mouvements parallèles, mais payer le taux fixe
du marché sur le notionnel a ramené le carry des titres à environ le taux variable.
\textbf{6.} 93~\%~: les dépôts non garantis subissent un retrait de 40~\%, et les MBS ne comptent qu'après la décote
et dans le plafond du niveau~2.
\textbf{7.} 36.1~milliards de dollars, soit 20.9~\% des dépôts.
\textbf{8.} SVB a perdu plus de 40~milliards de dollars le 9~mars, soit environ un quart des quelque 165~milliards
impliqués par les chiffres de sa direction (140~milliards représentant 85~\%)~; la banque stylisée n'aurait pu y faire face.
\textbf{9.} Repo sur les Treasuries et les MBS, guichet d'escompte contre collatéral valorisé au prix de marché diminué des
décotes, avances de la Federal Home Loan Bank~; chacun prend du temps et oblige à valoriser les titres.
\textbf{10.} Ce dispositif prêtait contre des titres au pair et non à la valeur de marché, pour un an au plus~: la banque pouvait emprunter
sur ses MBS sans réaliser la perte de 19~\%.
\textbf{11.} Les dépôts de particuliers sont garantis, nombreux et dispersés, et bougent lentement~; les dépôts d'entreprises non garantis
sont importants, concentrés et informés, et bougent ensemble.
\textbf{12.} La NII baisse lorsque le coût des dépôts augmente, alors que les actifs à taux fixe ne se réajustent pas~; l'EVE baisse
davantage, parce que les dépôts valent moins comme source de financement bon marché.
\textbf{13.} Le fait que la banque ait dû vendre à perte et ait eu besoin de capital~: un signal de faiblesse qui a incité
à la panique.
\textbf{14.} Parce que la panique porte sur la liquidité à valeur de marché~: les déposants ne peuvent être remboursés que
à partir d'actifs dont la valeur de marché est inférieure à leur valeur comptable, et chacun craint d'être le dernier servi.
\textbf{15.} Aucun des ratios usuels~: la concentration et les liens entre déposants demandent leurs propres
indicateurs (part des dépôts non garantis, plus gros déposants, concentration sectorielle).
\textbf{16.} Couvrir la duration des titres, les raccourcir, fixer le taux des dépôts avec un bêta réaliste,
et détenir de la liquidité pour la base des dépôts non garantis.
\textbf{17.} Au moins dans les mesures prudentielles et les \omterm{def:rc:stress-testing-and-scenarios:stress}{tests de résistance}~: un capital qui les ignore surestime
la capacité d'absorption des pertes en cas de panique.
\textbf{18.} Face à l'historique par cohorte et par cycle de taux, face aux pairs, avec des scénarios de tension de
sorties rapides, et en examinant la concentration de la base.
\textbf{19.} À retenir, \emph{mars 2023}~: la perte latente HTM de la banque stylisée représente 108~\% de ses
fonds propres, et ses actifs liquides s'épuisent après une sortie d'un jour de 20.9~\% des dépôts (36.1~milliards de dollars).
\textbf{20.} Les deux~: le risque de marché a rendu la banque insolvable à la valeur de marché, et le risque de liquidité
a transformé cette insolvabilité en faillite en une journée.
\end{solution}

\begin{solution}{iq:rc:liquidity-and-funding-risk-bank-treasury:1}
LCR~: HQLA face aux sorties nettes en situation de tension sur 30 jours, un coussin de survie à court terme. NSFR~: financement stable
face aux besoins de financement stable des actifs sur un an, une limite à la transformation d'échéances.
\iqlookfor{les deux horizons et ce que chacun contraint.}
\end{solution}

\begin{solution}{iq:rc:liquidity-and-funding-risk-bank-treasury:2}
Un prix interne des fonds par échéance (courbe de base plus prime de liquidité), facturé aux activités d'actif
et crédité aux activités de passif~; il mesure les marges nettes de financement, centralise le risque de taux et
de liquidité auprès de la trésorerie, et tarife l'usage de la liquidité.
\iqlookfor{le mécanisme et les incitations.}
\end{solution}

\begin{solution}{iq:rc:liquidity-and-funding-risk-bank-treasury:3}
L'EVE valorise l'ensemble des flux futurs~; la NII mesure l'année suivante. Une banque dont les actifs à taux fixe sont longs et les
dépôts stables perd de l'EVE lorsque les taux montent (ses actifs baissent), mais gagne de la NII pendant un temps (les actifs à taux variable
se réajustent, pas les dépôts).
\iqlookfor{différence d'horizon et un exemple.}
\end{solution}

\begin{solution}{iq:rc:liquidity-and-funding-risk-bank-treasury:4}
Les soldes se répartissent entre dépôts de cœur et dépôts volatils, avec des taux de décroissance estimés sur l'historique par segment~; les taux de
dépôt en fonction des taux de marché (bêta, décalages, planchers)~; surcouches de tension pour les sorties~; validation par
cycle, y compris lors des hausses de taux.
\iqlookfor{soldes, taux et validation.}
\end{solution}

\begin{solution}{iq:rc:liquidity-and-funding-risk-bank-treasury:5}
Titres de longue duration financés par des dépôts concentrés et non garantis, risque de taux non couvert, recours à la
comptabilisation des titres détenus jusqu'à l'échéance, \omterm{def:rc:stress-testing-and-scenarios:stress}{tests de résistance} de liquidité faibles, et une annonce publique qui a déclenché une
panique coordonnée.
\iqlookfor{duration, concentration du financement et gouvernance.}
\end{solution}

\begin{solution}{iq:rc:liquidity-and-funding-risk-bank-treasury:6}
Swaps payeurs de taux fixe adossés aux titres, ou receveurs de taux variable sur la duration modélisée des dépôts~; cela
réduit la NII lorsque les courbes sont à pente ascendante, et nécessite une comptabilité de couverture pour éviter la volatilité des
résultats publiés.
\iqlookfor{instrument, coût du carry et comptabilité.}
\end{solution}
